\chapter{Lie groups}
\label{chap:liegroups}

Throughout, $E$ and $F$ denote real Hausdorff locally convex spaces, not assumed complete.

\section{The definition}

\begin{definition}[Lie groups]
  \label{def:lie-group}
  \lean{LocallyConvex.LieGroups.MBLieGroup}
  \leanok
  A \emph{Lie group} modeled on $E$ is a manifold $G$ modeled on $E$, endowed with a group
  structure such that the multiplication $m_G\colon G \times G \to G$ and the inversion
  $\eta_G\colon G \to G$ are smooth \cite[Definition 5.1.1(b)]{GN26}.
\end{definition}

\begin{proposition}[Translations and conjugations]
  \label{prop:translations}
  \lean{LocallyConvex.LieGroups.MBLieGroup.smoothMBM_mul_left,
    LocallyConvex.LieGroups.MBLieGroup.smoothMBM_mul_right,
    LocallyConvex.LieGroups.MBLieGroup.smoothMBM_conj}
  \leanok
  \uses{def:lie-group}
  For $g \in G$, the left multiplication $\lambda_g(x) = gx$, the right multiplication
  $\rho_g(x) = xg$ and the conjugation $x \mapsto gxg^{-1}$ are smooth.
\end{proposition}

\begin{proof}
  \leanok
  $\lambda_g = m_G \circ (x \mapsto (g, x))$, and a pair of smooth maps is smooth; similarly for
  $\rho_g$, and the conjugation is $\rho_{g^{-1}} \circ \lambda_g$.
\end{proof}

\begin{proposition}[Lie groups are topological groups]
  \label{prop:topological-group}
  \lean{LocallyConvex.LieGroups.MBLieGroup.isTopologicalGroup}
  \leanok
  \uses{def:lie-group}
  Every Lie group is a topological group.
\end{proposition}

\begin{proof}
  \leanok
  Smooth maps are continuous, and the topology of the manifold $G \times G$ is the product
  topology.
\end{proof}

\begin{definition}[Morphisms]
  \label{def:lie-hom}
  \lean{LocallyConvex.LieGroups.MBLieGroup.IsHom}
  \leanok
  \uses{def:lie-group}
  A \emph{morphism of Lie groups} is a smooth group homomorphism \cite[Definition
  5.1.1(c)]{GN26}.
\end{definition}

\begin{proposition}
  \label{prop:lie-hom-cat}
  \lean{LocallyConvex.LieGroups.MBLieGroup.IsHom.id, LocallyConvex.LieGroups.MBLieGroup.IsHom.comp}
  \leanok
  \uses{def:lie-hom}
  Identities are morphisms of Lie groups, and composites of morphisms are morphisms.
\end{proposition}

\begin{proof}
  \leanok
  The chain rule.
\end{proof}

\section{The group of diffeomorphisms of a compact manifold}

Let $K$ be a compact manifold with a local addition. The diffeomorphisms of $K$ form an open
subset $\mathrm{Diff}(K)$ of the manifold $C^\infty(K, K)$ \cite[Proposition 13.1.10]{GN26}
(proved in \cite{LCM}), hence a manifold.

\begin{theorem}[Smoothness of the inversion]
  \label{thm:diff-inv}
  \lean{LocallyConvex.LieGroups.smooth_inv, LocallyConvex.LieGroups.diffeoAtlas}
  \leanok
  The inversion $\mathrm{Diff}(K) \to C^\infty(K, K)$, $\gamma \mapsto \gamma^{-1}$, is smooth
  \cite[Proposition 13.1.10(a)]{GN26}.
\end{theorem}

\begin{proof}
  \leanok
  By the exponential law it suffices that $(\gamma, y) \mapsto \gamma^{-1}(y)$ is $C^{k,\ell}$ on
  $\mathrm{Diff}(K) \times K$ for all $k$, $\ell$. In charts, $\gamma^{-1}(y)$ solves the equation
  $c(\gamma(e^{-1}(z))) = c(y)$ for $z$, with parameters $(\gamma, y)$ in a locally convex space
  and the unknown $z$ in the finite-dimensional model of $K$; the solution is continuous, since
  its graph is closed and $K$ is compact. The implicit function theorem with locally convex
  parameters and finite-dimensional values \cite{LCA} then gives smoothness
  \cite[Lemma 13.1.8]{GN26}.
\end{proof}
